\documentclass[11pt]{book}
\usepackage[paperwidth=145mm,paperheight=200mm,top=18mm,bottom=20mm,inner=20mm,outer=16mm]{geometry}
\usepackage{brumfiel}
% figures07.tex -- Chapter 7 figures redrawn in TikZ.
% Geometry reconstructed from sources/Geometry.pdf pp.130-146 (book pp.116-132)
% and checked against scans/scan08.pdf pp.11-24 and scans/scan09.pdf pp.1-6.
% Every constrained point is computed, never eyeballed; see tools/constraints/ch07.py.

% ---------------------------------------------------------------- Fig 7-1
\newcommand{\FIGVIIONE}{\begin{bkfigure}[0.60\linewidth]
\begin{tikzpicture}[scale=1.235]
  \coordinate (A) at (0,0); \coordinate (B) at (3.1,0);
  \coordinate (D) at (3.9,0); \coordinate (C) at (1.64,2.1);
  \draw[given] (A) -- (D);
  \draw[given] (A) -- (C) -- (B);
  \foreach \p/\d in {A/below left, B/below, D/below right, C/above}
    {\node[vlab,\d] at (\p) {$\p$};}
\end{tikzpicture}
\figcap{7--1}
\end{bkfigure}}

% ------------------------------------------------------------ Fig 7-2, 7-3
\newcommand{\FIGVIITWOTHREE}{\begin{bkfigure}[\linewidth]
\begin{minipage}{0.48\linewidth}\centering
\begin{tikzpicture}[scale=0.82]
  % BD is chosen congruent to AC (book's own wording in the First Proof),
  % so D sits at distance |AC| beyond B.  |AC| = sqrt(1.10^2+2.10^2).
  \coordinate (A) at (0,0); \coordinate (B) at (2.6,0);
  \coordinate (C) at (1.10,2.10);
  \coordinate (D) at (4.97066,0);
  \draw[given] (A) -- (5.28,0);
  \draw[given] (A) -- (C) -- (B);
  \draw[aux] (C) -- (D);
  % angle ACB at C, angle DBC at B -- the pair the proof assumes congruent
  \draw[tick] ($(C)+(242.35:0.58)$) arc (242.35:305.54:0.58);
  \draw[tick] ($(B)+(0:0.48)$) arc (0:125.54:0.48);
  \foreach \p/\d in {A/below left, B/below, D/below right, C/above}
    {\node[vlab,\d] at (\p) {$\p$};}
\end{tikzpicture}
\figcap{7--2}
\end{minipage}\hfill
\begin{minipage}{0.48\linewidth}\centering
\begin{tikzpicture}[scale=0.82]
  \coordinate (A) at (0,0); \coordinate (B) at (2.4,0);
  \coordinate (D) at (4.29,0); \coordinate (C) at (1.2,2.1);
  \coordinate (F) at (1.05,0);
  \coordinate (E) at ($(C)!1.245!(F)$);
  \draw[given] (A) -- (D);
  \draw[given] (A) -- (C) -- (B);
  \draw[aux] (C) -- (E);
  % angle ECB at C -- the angle the proof makes congruent to angle DBC
  \draw[tick] ($(C)+(265.91:0.72)$) arc (265.91:299.74:0.72);
  \foreach \p/\d in {A/below left, B/below, D/below right, C/above,
                     F/above right, E/below}
    {\node[vlab,\d] at (\p) {$\p$};}
\end{tikzpicture}
\figcap{7--3}
\end{minipage}
\end{bkfigure}}

% ---------------------------------------------------------------- Fig 7-4
\newcommand{\FIGVIIFOUR}{\begin{bkfigure}[0.62\linewidth]
\begin{tikzpicture}[scale=1.235]
  \coordinate (A) at (0,0); \coordinate (B) at (2.6,0);
  \coordinate (D) at (5.05,0); \coordinate (C) at (1.15,2.1);
  \coordinate (G) at ($(C)!0.5!(B)$);
  \coordinate (H) at ($(A)!2!(G)$);
  \draw[given] (A) -- (D);
  \draw[given] (A) -- (C) -- (B);
  \draw[given] (B) -- (H);
  \draw[aux] (A) -- (H);
  \foreach \p/\d in {A/below left, B/below, D/below right, C/above,
                     H/above right}
    {\node[vlab,\d] at (\p) {$\p$};}
  \node[vlab] at (1.957,1.40) {$G$};
\end{tikzpicture}
\figcap{7--4}
\end{bkfigure}}

% ---------------------------------------------------------------- Fig 7-5
\newcommand{\FIGVIIFIVE}{\begin{bkfigure}[0.72\linewidth]
\begin{tikzpicture}[scale=1.235]
  \coordinate (P) at (1.30,1.95);
  \coordinate (Q) at (1.15,0); \coordinate (R) at (1.75,0);
  \draw[given] (0,0) -- (3.2,0);
  \draw[given] (P) -- (Q); \draw[given] (P) -- (R);
  \draw[tick] (1.15,0.16) -- (1.31,0.16) -- (1.31,0);
  \draw[tick] (1.75,0.16) -- (1.59,0.16) -- (1.59,0);
  \node[vlab,above] at (P) {$P$};
  \node[vlab,below] at (Q) {$Q$};
  \node[vlab,below] at (R) {$R$};
  \node[vlab,right] at (3.2,0) {$l$};
\end{tikzpicture}
\figcap{7--5}
\end{bkfigure}}

% ---------------------------------------------------------------- Fig 7-6
\newcommand{\FIGVIISIX}{\begin{bkfigure}[0.58\linewidth]
\begin{tikzpicture}[scale=1.235]
  \coordinate (B) at (0,0); \coordinate (C) at (2.7,0);
  \coordinate (A) at (2.5,1.9);
  % C' on AB with AC' = AC  (t = |AC|/|AB|)
  \coordinate (Cp) at ($(A)!0.60843!(B)$);
  \draw[given] (B) -- (C) -- (A) -- cycle;
  \draw[aux] (Cp) -- (C);
  \node[vlab,above] at (A) {$A$};
  \node[vlab,below left] at (B) {$B$};
  \node[vlab,below right=2pt] at (C) {$C$};
  \node[vlab] at (0.737,1.062) {$C'$};
\end{tikzpicture}
\figcap{7--6}
\end{bkfigure}}

% ------------------------------------------------------------ Fig 7-7, 7-8
\newcommand{\FIGVIISEVENEIGHT}{\begin{bkfigure}[\linewidth]
\begin{minipage}{0.48\linewidth}\centering
\begin{tikzpicture}[scale=1.197]
  % sides 17, 16, 25 to scale (0.1 unit per unit length)
  \coordinate (A) at (0,0); \coordinate (C) at (2.5,0);
  \coordinate (B) at (1.316,1.076);
  \draw[given] (A) -- (C) -- (B) -- cycle;
  \node[vlab,below left] at (A) {$A$};
  \node[vlab,below right] at (C) {$C$};
  \node[vlab,above] at (B) {$B$};
  \node[slab,above left] at ($(A)!0.5!(B)$) {17};
  \node[slab,above right] at ($(B)!0.5!(C)$) {16};
  \node[slab,below] at ($(A)!0.5!(C)$) {25};
\end{tikzpicture}
\figcap{7--7}
\end{minipage}\hfill
\begin{minipage}{0.48\linewidth}\centering
\begin{tikzpicture}[scale=1.197]
  \coordinate (P) at (0,0); \coordinate (Q) at (2.6,0.35);
  \coordinate (R) at (1.10,2.55);
  % QS bisects angle PQR  =>  PS/SR = QP/QR
  \coordinate (S) at ($(P)!0.49629!(R)$);
  \draw[given] (P) -- (Q) -- (R) -- cycle;
  \draw[given] (S) -- (Q);
  % congruence ticks across the two halves of angle PQR
  \draw[tick] ([shift=(65.95:0.09)]$(Q)!0.30!(S)$)
            -- ([shift=(245.95:0.09)]$(Q)!0.30!(S)$);
  \draw[tick] ([shift=(34.30:0.09)]$(Q)!0.30!(R)$)
            -- ([shift=(214.30:0.09)]$(Q)!0.30!(R)$);
  \node[vlab,below left] at (P) {$P$};
  \node[vlab,right] at (Q) {$Q$};
  \node[vlab,above] at (R) {$R$};
  \node[vlab,above left] at (S) {$S$};
\end{tikzpicture}
\figcap{7--8}
\end{minipage}
\end{bkfigure}}

% ----------------------------------------------------------- Fig 7-9, 7-10
\newcommand{\FIGVIININETEN}{\begin{bkfigure}[\linewidth]
\begin{minipage}{0.46\linewidth}\centering
\begin{tikzpicture}[scale=1.175]
  \coordinate (A) at (0,0); \coordinate (C) at (2.9,0);
  \coordinate (B) at (1.0,1.62);
  % B' on AC with AB' = AB   (|AB| = sqrt(1.0^2 + 1.62^2))
  \coordinate (Bp) at (1.903786,0);
  \draw[given] (A) -- (C) -- (B) -- cycle;
  \draw[aux] (B) -- (Bp);
  % the book's two angle arcs: one at B crossed by BB' (angles 1, 4),
  % one at B' crossed by B'B (angles 2, 3)
  \draw[tick] ($(B)+(238.33:0.42)$) arc (238.33:319.54:0.42);
  \draw[tick] ($(Bp)+(0:0.20)$) arc (0:180:0.20);
  \node[vlab,below left] at (A) {$A$};
  \node[vlab,below right] at (C) {$C$};
  \node[vlab,above] at (B) {$B$};
  \node[vlab,below] at (Bp) {$B'$};
  \node[slab] at (0.986,1.000) {1};   % angle ABB'
  \node[slab] at (1.665,0.808) {4};   % angle B'BC
  \node[slab] at (1.429,0.278) {2};   % angle AB'B
  \node[slab] at (2.210,0.300) {3};   % angle BB'C
\end{tikzpicture}
\figcap{7--9}
\end{minipage}\hfill
\begin{minipage}{0.50\linewidth}\centering
\begin{tikzpicture}[scale=0.729]
  % vertex ratios measured off scan08 p.15 (grid overlay); three interior
  % vertices are unlabelled in the book, as here.
  \coordinate (Pone) at (0.73,1.58);
  \coordinate (Ptwo) at (1.61,4.41);
  \coordinate (Pthree) at (1.73,1.61);
  \coordinate (U) at (0.12,2.99);
  \coordinate (R) at (2.78,2.16);
  \coordinate (S) at (2.38,0.83);
  \coordinate (Pn1) at (0.34,0.45);
  \coordinate (Pn) at (2.60,0.10);
  \draw[given] (Pone) -- (Ptwo) -- (Pthree) -- (U) -- (R) -- (S)
               -- (Pn1) -- (Pn);
  \draw[aux] (Pone) -- (Pn);
  \node[vlab,left] at (Pone) {$P_1$};
  \node[vlab,above] at (Ptwo) {$P_2$};
  \node[vlab,right] at (Pthree) {$P_3$};
  \node[vlab,left] at (Pn1) {$P_{n-1}$};
  \node[vlab,right] at (Pn) {$P_n$};
\end{tikzpicture}
\figcap{7--10}
\end{minipage}
\end{bkfigure}}

% --------------------------------------------------------------- Fig 7-11
\newcommand{\FIGVIIELEVEN}{\begin{bkfigure}[0.56\linewidth]
\begin{tikzpicture}[scale=1.300]
  \coordinate (B) at (0,0); \coordinate (C) at (2.6,0);
  \coordinate (A) at (1.85,1.35);
  % D on BA with DB = DC  (D lies on the perpendicular bisector x = 1.3)
  \coordinate (D) at (1.3,0.94865);
  \draw[given] (B) -- (C) -- (A) -- cycle;
  \draw[given] (D) -- (C);
  % DB = DC, so the book marks angle DBC and angle DCB congruent
  % with a DOUBLE arc at each vertex.
  \draw[tick] (0.30,0) arc (0:36.13:0.30);
  \draw[tick] (0.40,0) arc (0:36.13:0.40);
  \draw[tick] ($(C)+(180:0.30)$) arc (180:143.87:0.30);
  \draw[tick] ($(C)+(180:0.40)$) arc (180:143.87:0.40);
  \node[vlab,below left] at (B) {$B$};
  \node[vlab,below right] at (C) {$C$};
  \node[vlab,above] at (A) {$A$};
  \node[vlab,above left] at (D) {$D$};
\end{tikzpicture}
\figcap{7--11}
\end{bkfigure}}

% --------------------------------------------------------------- Fig 7-12
\newcommand{\FIGVIITWELVE}{\begin{bkfigure}[0.62\linewidth]
\begin{tikzpicture}[scale=1.300]
  % C is on the perpendicular bisector of BB', so CB' = CB exactly
  % (the hypothesis BC = B'C' with C' = C).  Q on AB with QB = QB'.
  \coordinate (A) at (0,0); \coordinate (B) at (3.2,0);
  \coordinate (C) at (1.8875,1.43);
  \coordinate (Bp) at (1.30,-0.42);     % CB' = CB
  \coordinate (Q) at (2.2036,0);        % QB = QB'
  \draw[given] (A) -- (B);
  \draw[given] (A) -- (C) -- (B);
  \draw[given] (C) -- (Bp);
  \draw[aux] (Bp) -- (Q);
  \node[vlab,above] at (C) {$C'=C$};
  \node[vlab,left] at (A) {$A'=A$};
  \node[vlab,left] at (Bp) {$B'$};
  \node[vlab,above right] at (Q) {$Q$};
  \node[vlab,right] at (B) {$B$};
\end{tikzpicture}
\figcap{7--12}
\end{bkfigure}}

% ------------------------------------------------------- Fig 7-13, 7-14
\newcommand{\FIGVIITHIRTEENFOURTEEN}{\begin{bkfigure}[\linewidth]
\begin{minipage}{0.48\linewidth}\centering
\begin{tikzpicture}[scale=1.145]
  \coordinate (A) at (0,0); \coordinate (C) at (3.0,0);
  \coordinate (B) at (1.45,1.93);
  \coordinate (E) at ($(B)!0.49!(C)$);
  \coordinate (D) at ($(A)!0.655!(E)$);
  \draw[given] (A) -- (C) -- (B) -- cycle;
  \draw[given] (A) -- (D);
  \draw[aux] (D) -- (E);
  \draw[given] (D) -- (C);
  \node[vlab,below left] at (A) {$A$};
  \node[vlab,below right] at (C) {$C$};
  \node[vlab,above] at (B) {$B$};
  \node[vlab,above] at (D) {$D$};
  \node[vlab,above right] at (E) {$E$};
\end{tikzpicture}
\figcap{7--13}
\end{minipage}\hfill
\begin{minipage}{0.48\linewidth}\centering
\begin{tikzpicture}[scale=1.145]
  \draw[given] (1.2,2.03) -- (1.2,-0.10);
  \draw[given] (0,1.6) -- (2.8,1.6);
  \draw[given] (0,0.2) -- (2.8,0.2);
  % both right-angle boxes sit ABOVE the horizontal, right of m (as printed)
  \draw[tick] (1.2,1.76) -- (1.36,1.76) -- (1.36,1.6);
  \draw[tick] (1.2,0.36) -- (1.36,0.36) -- (1.36,0.2);
  \node[vlab,above] at (1.2,2.03) {$m$};
  \node[vlab,right] at (2.8,1.6) {$l$};
  \node[vlab,right] at (2.8,0.2) {$l'$};
\end{tikzpicture}
\figcap{7--14}
\end{minipage}
\end{bkfigure}}

% --------------------------------------------------------------- Fig 7-15
\newcommand{\FIGVIIFIFTEEN}{\begin{bkfigure}[\linewidth]
\begin{minipage}{0.48\linewidth}\centering
\begin{tikzpicture}[scale=1.184]
  \def\rr{1.15}
  \coordinate (O) at (0,0);
  \draw[given] (O) circle (\rr);
  \coordinate (A) at (40:\rr);          % on the circle
  \coordinate (Bi) at (95:0.80);        % interior: OB < r
  \coordinate (Cx) at (-10:1.87);       % exterior: OC > r
  \coordinate (E) at (215:\rr);
  \coordinate (F) at (293:\rr);
  \draw[given] (O) -- (A);
  \draw[given] (O) -- (Bi);
  \draw[given] (O) -- (Cx);
  \draw[given] (E) -- (F);            % the chord EF
  \draw[aux] (E) -- ($(F)!-0.30!(E)$);  % secant, produced both ways
  \draw[aux] (F) -- ($(E)!-0.30!(F)$);
  \node[vlab,above right] at (A) {$A$};
  \node[vlab,left] at (Bi) {$B$};
  \node[vlab,right] at (Cx) {$C$};
  \node[vlab,below left] at (E) {$E$};
  \node[vlab,below] at (F) {$F$};
  \node[vlab,below left] at (O) {$O$};
\end{tikzpicture}
\end{minipage}\hfill
\begin{minipage}{0.48\linewidth}\centering
\begin{tikzpicture}[scale=1.184]
  \def\rr{1.15}
  \coordinate (O) at (0,0);
  \draw[given] (O) circle (\rr);
  \coordinate (A) at (95:\rr);
  \coordinate (B) at (45:\rr);
  \coordinate (F) at (178:\rr);
  \coordinate (E) at (250:\rr);
  \coordinate (G) at (310:\rr);
  \draw[given] (O) -- (A);            % central angle AOB
  \draw[given] (O) -- (B);
  % inscribed angle EFG: only the two chords FE and FG are drawn.
  % (Checked on scan08 p.18 -- the book does NOT close E--G.)
  \draw[given] (F) -- (E);
  \draw[given] (F) -- (G);
  \node[vlab,above] at (A) {$A$};
  \node[vlab] at (38:1.48) {$B$};
  \node[vlab,left] at (F) {$F$};
  \node[vlab,below] at (E) {$E$};
  \node[vlab,below right] at (G) {$G$};
  \node[vlab,below right] at (O) {$O$};
\end{tikzpicture}
\end{minipage}
\figcap{7--15}
\end{bkfigure}}

% --------------------------------------------------------------- Fig 7-16
\newcommand{\FIGVIISIXTEEN}{\begin{bkfigure}[0.56\linewidth]
\begin{tikzpicture}[scale=1.365]
  % Ex. 7-5 #3 asks what kind of angle each of /A, /B, /C, /D, /E is, so
  % every labelled point must carry two drawn rays.  C is the CENTRE
  % (printed "C = O"), and the two radii CD, CP_3 form the central angle.
  % Chord set read off scan08 p.20; P_1..P_4 are the book's unlabelled
  % points where the other chords meet the circle.
  \def\rr{1.3}
  \coordinate (O) at (0,0);
  \draw[given] (O) circle (\rr);
  \coordinate (A) at (57:\rr);
  \coordinate (B) at (132:\rr);
  \coordinate (D) at (240:\rr);
  \coordinate (E) at (308:\rr);
  \coordinate (Pone) at (-4:\rr);
  \coordinate (Ptwo) at (13:\rr);
  \coordinate (Pthree) at (-13:\rr);
  \coordinate (Pfour) at (-34:\rr);
  \draw[given] (B) -- (D);                       % chord, side of /B and /D
  \draw[given] (B) -- (Pone);                    % chord, other side of /B
  \draw[given] (A) -- (Ptwo);                    % chord, one side of /A
  \draw[given] (A) -- ($(A)+(-19.5:1.05)$);      % ray outside, other side of /A
  \draw[given] (O) -- (D);                       % radius
  \draw[given] (O) -- (Pthree);                  % radius
  \draw[given] (D) -- (Pfour);                   % chord
  \draw[given] (D) -- (E);                       % chord
  \draw[given] (E) -- ($(E)+(9.9:1.20)$);        % the two rays at E, both
  \draw[given] (E) -- ($(E)+(-11.5:1.20)$);      % lying outside the circle
  \node[vlab,above left] at (B) {$B$};
  \node[vlab,above right] at (A) {$A$};
  \node[vlab,below left] at (D) {$D$};
  \node[vlab,below] at (E) {$E$};
  \node[slab] at (-0.30,0.22) {$C=O$};
\end{tikzpicture}
\figcap{7--16}
\end{bkfigure}}

% --------------------------------------------------------------- Fig 7-17
% Line drawing of a drafting compass, brace beneath marking the radius r.
\newcommand{\FIGVIISEVENTEEN}{\begin{bkfigure}[0.40\linewidth]
\begin{tikzpicture}[scale=1.300]
  % handle
  \draw[given] (-0.075,2.62) rectangle (0.075,3.05);
  % knurled collar
  \draw[given] (-0.115,2.30) -- (-0.115,2.62) -- (0.115,2.62) -- (0.115,2.30);
  \draw[given] (-0.115,2.30) arc (180:360:0.115);
  % hinge head
  \draw[given] (-0.16,1.92) .. controls (-0.20,2.16) and (-0.12,2.28) ..
               (0,2.30) .. controls (0.12,2.28) and (0.20,2.16) .. (0.16,1.92);
  % legs
  \draw[given] (-0.13,2.02) -- (-0.58,0.30);
  \draw[given] (0.13,2.02) -- (0.58,0.30);
  \draw[given] (-0.055,2.00) -- (-0.50,0.30);
  \draw[given] (0.055,2.00) -- (0.50,0.30);
  % needle point (left) and pen nib (right)
  \draw[given] (-0.58,0.30) -- (-0.605,0.10) -- (-0.545,0.10) -- (-0.50,0.30);
  \draw[given] (-0.605,0.10) -- (-0.62,0.0);
  \draw[given] (0.50,0.30) -- (0.545,0.06) -- (0.605,0.06) -- (0.58,0.30);
  \draw[given] (0.575,0.06) -- (0.62,0.0);
  % brace showing the radius
  \draw[tick] (-0.62,-0.16) -- (-0.62,-0.24) -- (-0.04,-0.24);
  \draw[tick] (0.0,-0.30) -- (0.04,-0.24) -- (0.62,-0.24) -- (0.62,-0.16);
  \node[vlab,below] at (0,-0.34) {$r$};
\end{tikzpicture}
\figcap{7--17}
\end{bkfigure}}

% ------------------------------------------------------- Fig 7-18, 7-19
\newcommand{\FIGVIIEIGHTEENNINETEEN}{\begin{bkfigure}[\linewidth]
\begin{minipage}{0.48\linewidth}\centering
\begin{tikzpicture}[scale=1.176]
  \coordinate (V) at (0,0);
  \foreach \a in {45,18,-9,-36}{\draw[given] (V) -- (\a:2.35);}
  \draw[tick] (-36:1.75) arc (-36:45:1.75);
  \node[slab] at (31.5:1.20) {$B$};
  \node[slab] at (4.5:1.20) {$C$};
  \node[slab] at (-22.5:1.20) {$D$};
  \node[slab] at (2.020,0.324) {$A$};
\end{tikzpicture}
\figcap{7--18}
\end{minipage}\hfill
\begin{minipage}{0.48\linewidth}\centering
\begin{tikzpicture}[scale=1.176]
  \draw[given] (0,0) circle (0.85);
  \draw[given] (1.6,-0.85) rectangle (3.3,0.85);
\end{tikzpicture}
\figcap{7--19}
\end{minipage}
\end{bkfigure}}

% ------------------------------------------------------- Fig 7-20, 7-21
\newcommand{\FIGVIITWENTYTWENTYONE}{\begin{bkfigure}[\linewidth]
\begin{minipage}{0.48\linewidth}\centering
\begin{tikzpicture}[scale=0.998]
  % angle A
  \coordinate (A) at (0,1.75);
  \draw[given] (A) -- ++(0:3.0);
  \draw[given] (A) -- ++(14:3.0);
  \node[vlab,left] at (A) {$A$};
  % ray from B, with the two possible copied angles dotted in
  \coordinate (B) at (0,0);
  \draw[given] (B) -- ++(0:3.0);
  \draw[aux] (B) -- ++(20:2.5);
  \draw[aux] (B) -- ++(-20:2.5);
  \node[vlab,left] at (B) {$B$};
  \node[slab,right] at (3.0,0) {Ray};
\end{tikzpicture}
\figcap{7--20}
\end{minipage}\hfill
\begin{minipage}{0.48\linewidth}\centering
\begin{tikzpicture}[scale=0.998]
  % copied angle at A
  \coordinate (A) at (0,1.75);
  \coordinate (E) at ($(A)+(0:1.0)$);
  \coordinate (D) at ($(A)+(38:1.0)$);
  \draw[given] (A) -- ++(0:2.9);
  \draw[given] (A) -- ++(38:1.85);
  % the compass arc of radius r about A, cutting both sides (at D and E)
  \draw[aux] ($(A)+(-8:1.0)$) arc (-8:50:1.0);
  \node[vlab,left] at (A) {$A$};
  \node[vlab] at ($(A)+(50:1.35)$) {$D$};
  \node[vlab,below right] at (E) {$E$};
  % construction at B
  \coordinate (B) at (0,0);
  \coordinate (F) at ($(B)+(0:1.0)$);
  \coordinate (G) at ($(B)+(38:1.0)$);
  \draw[given] (B) -- ++(0:2.9);
  \draw[given] (B) -- ($(B)+(38:1.85)$);
  % arc of radius r about B, cutting the ray at F;
  % then the arc of radius DE about F, which meets the first at G.
  \draw[aux] ($(B)+(-8:1.0)$) arc (-8:50:1.0);
  \draw[aux] ($(F)+(96:0.6512)$) arc (96:122:0.6512);
  \node[vlab,left] at (B) {$B$};
  \node[vlab,below] at (F) {$F$};
  \node[vlab] at ($(B)+(50:1.35)$) {$G$};
  \node[vlab,above right] at ($(B)+(38:1.85)$) {$H$};
  \node[slab,right] at (2.9,0) {Ray};
\end{tikzpicture}
\figcap{7--21}
\end{minipage}
\end{bkfigure}}

% ------------------------------------------------------- Fig 7-22, 7-23
\newcommand{\FIGVIITWENTYTWOTWENTYTHREE}{\begin{bkfigure}[\linewidth]
\begin{minipage}{0.48\linewidth}\centering
\begin{tikzpicture}[scale=0.983]
  % given triangle
  \coordinate (A) at (0,2.05); \coordinate (B) at (0.95,3.40);
  \coordinate (C) at (2.35,2.05);
  \draw[given] (A) -- (B) -- (C) -- cycle;
  \node[vlab,left] at (A) {$A$};
  \node[vlab,above] at (B) {$B$};
  \node[vlab,right] at (C) {$C$};
  % copy on a ray
  \coordinate (Ap) at (0,0); \coordinate (Bp) at (0.95,1.35);
  \coordinate (Cp) at (2.35,0);   % A'C' = AC, A'B' = AB (congruent copy)
  \draw[given] (Ap) -- (3.05,0);
  % Construction 7-2 stops here: the copied side A'B', the ray, C' cut off
  % on it, and the arc of radius AB about A' that fixes B'.  The book does
  % NOT close B'C' (scan08 p.24).
  \draw[given] (Ap) -- (Bp);
  \draw[aux] ($(Ap)+(44.9:1.65076)$) arc (44.9:64.9:1.65076);
  \draw[tick] ($(Cp)+(0,-0.13)$) -- ($(Cp)+(0,0.13)$);
  \node[vlab,left] at (Ap) {$A'$};
  \node[vlab] at ($(Bp)+(0.10,0.36)$) {$B'$};
  \node[vlab] at ($(Cp)+(0.26,0.34)$) {$C'$};
\end{tikzpicture}
\figcap{7--22}
\end{minipage}\hfill
\begin{minipage}{0.48\linewidth}\centering
\begin{tikzpicture}[scale=0.983]
  % given triangle
  \coordinate (A) at (0,2.15); \coordinate (B) at (1.15,3.45);
  \coordinate (C) at (2.4,2.15);
  \draw[given] (A) -- (B) -- (C) -- cycle;
  \node[vlab,left] at (A) {$A$};
  \node[vlab,above] at (B) {$B$};
  \node[vlab,right] at (C) {$C$};
  % ray with the two arcs crossing at B'
  \coordinate (Ap) at (0,0); \coordinate (Cp) at (2.4,0);
  \coordinate (Bp) at (1.15,1.30);
  \draw[given] (Ap) -- (3.1,0);
  \dt{Ap}
  \draw[tick] ($(Cp)+(0,-0.13)$) -- ($(Cp)+(0,0.13)$);
  % the two construction circles: radius AB about A', radius BC about C'.
  % They cross at B'.
  \draw[aux] ($(Ap)+(38.5:1.73566)$) arc (38.5:58.5:1.73566);
  \draw[aux] ($(Cp)+(123.9:1.80347)$) arc (123.9:143.9:1.80347);
  \node[vlab,left] at (Ap) {$A'$};
  \node[vlab] at ($(Bp)+(0,0.44)$) {$B'$};
  \node[vlab] at ($(Cp)+(0.26,0.34)$) {$C'$};
  \node[slab,right] at (3.1,0) {Ray};
\end{tikzpicture}
\figcap{7--23}
\end{minipage}
\end{bkfigure}}

% ------------------------------------------------------- Fig 7-24, 7-25
\newcommand{\FIGVIITWENTYFOURTWENTYFIVE}{\begin{bkfigure}[\linewidth]
\begin{minipage}{0.48\linewidth}\centering
\begin{tikzpicture}[scale=0.993]
  \coordinate (A) at (0,0); \coordinate (B) at (3.0,0);
  \coordinate (P) at (1.5,1.0828); \coordinate (Q) at (1.5,-1.0828);
  \draw[given] (A) -- (B);
  \dt{A} \dt{B}
  \draw[aux] ([shift=(-50:1.85)]A) arc (-50:50:1.85);
  \draw[aux] ([shift=(130:1.85)]B) arc (130:230:1.85);
  \node[vlab,left] at (A) {$A$};
  \node[vlab,right] at (B) {$B$};
  \node[vlab] at (2.35,1.30) {$P$};
  \node[vlab] at (2.35,-1.30) {$Q$};
\end{tikzpicture}
\figcap{7--24}
\end{minipage}\hfill
\begin{minipage}{0.48\linewidth}\centering
\begin{tikzpicture}[scale=0.993]
  \coordinate (A) at (0,0); \coordinate (B) at (3.0,0);
  \coordinate (P) at (1.5,1.0828); \coordinate (Q) at (1.5,-1.0828);
  \coordinate (C) at (1.5,0);
  \draw[given] (A) -- (B);
  \dt{A} \dt{B}
  \draw[aux] ([shift=(-50:1.85)]A) arc (-50:50:1.85);
  \draw[aux] ([shift=(130:1.85)]B) arc (130:230:1.85);
  \draw[given] ($(P)!-0.16!(Q)$) -- ($(Q)!-0.16!(P)$);
  \node[vlab,left] at (A) {$A$};
  \node[vlab,right] at (B) {$B$};
  \node[vlab] at (2.35,1.30) {$P$};
  \node[vlab] at (2.35,-1.30) {$Q$};
  \node[vlab] at (2.05,0.72) {$C$};
\end{tikzpicture}
\figcap{7--25}
\end{minipage}
\end{bkfigure}}

% -------------------------------------------------- Fig 7-26, 7-27, 7-28
\newcommand{\FIGVIITWENTYSIXSEVENEIGHT}{\begin{bkfigure}[\linewidth]
\begin{minipage}{0.30\linewidth}\centering
\begin{tikzpicture}[scale=0.80]
  \draw[given] (0,0) -- (3.0,0);
  \coordinate (A) at (0.55,0); \coordinate (P) at (1.5,0);
  \coordinate (B) at (2.45,0);
  \draw[tick] (0.55,-0.16) -- (0.55,0.16);
  \draw[tick] (2.45,-0.16) -- (2.45,0.16);
  \dt{P}
  \node[vlab,above left] at (A) {$A$};
  \node[vlab,above] at (P) {$P$};
  \node[vlab,above right] at (B) {$B$};
\end{tikzpicture}
\figcap{7--26}
\end{minipage}\hfill
\begin{minipage}{0.30\linewidth}\centering
\begin{tikzpicture}[scale=0.80]
  \draw[given] (0,0) -- (3.2,0);
  \coordinate (A) at (0.35,0); \coordinate (B) at (2.85,0);
  \coordinate (P) at (1.6,1.75);
  \draw[tick] (0.35,-0.16) -- (0.35,0.16);
  \draw[tick] (2.85,-0.16) -- (2.85,0.16);
  \draw[aux] (P) -- (1.6,0);
  \node[vlab,above] at (P) {$P$};
  \node[vlab,above right] at (A) {$A$};
  \node[vlab,above left] at (B) {$B$};
\end{tikzpicture}
\figcap{7--27}
\end{minipage}\hfill
\begin{minipage}{0.34\linewidth}\centering
\begin{tikzpicture}[scale=0.80]
  \coordinate (A) at (0,0);
  \coordinate (P) at (1.55,1.8); \coordinate (Pp) at (1.55,-1.8);
  \coordinate (D) at (1.55,0);
  \coordinate (B) at (2.45,0); \coordinate (C) at (2.9,0);
  \draw[given] (-0.4,0) -- (3.4,0);
  \draw[given] (A) -- (P);
  \draw[given] (A) -- ($(A)!1.22!(Pp)$);   % the ray AQ, through P'
  \draw[given] (P) -- (Pp);
  % angle PAC congruent to angle QAC -- the book marks both
  \draw[tick] (0.42,0) arc (0:49.28:0.42);
  \draw[tick] (0.42,0) arc (0:-49.28:0.42);
  \dt{B} \dt{C}
  \node[vlab] at (-0.24,0.46) {$A$};
  \node[vlab,above] at (P) {$P$};
  \node[vlab,below left] at (Pp) {$P'$};
  \node[vlab] at (1.80,0.30) {$D$};
  \node[vlab,above] at (B) {$B$};
  \node[vlab,above] at (C) {$C$};
  \node[vlab,right] at (3.4,0) {$l$};
  \node[vlab,below right] at ($(A)!1.22!(Pp)$) {$Q$};
\end{tikzpicture}
\figcap{7--28}
\end{minipage}
\end{bkfigure}}

% ------------------------------------------------------- Fig 7-29, 7-30
\newcommand{\FIGVIITWENTYNINETHIRTY}{\begin{bkfigure}[\linewidth]
\begin{minipage}{0.48\linewidth}\centering
\begin{tikzpicture}[scale=1.074]
  \coordinate (A) at (0,0);
  \coordinate (P) at (1.6,1.6); \coordinate (Pp) at (1.6,-1.6);
  \draw[given] (-0.3,0) -- (3.4,0);
  \draw[given] (A) -- (P); \draw[given] (A) -- (Pp);
  \draw[aux] (P) -- (Pp);
  \node[vlab] at (-0.18,0.26) {$A$};
  \node[vlab,above] at (P) {$P$};
  \node[vlab,below] at (Pp) {$P'$};
  \node[vlab,right] at (3.4,0) {$l$};
  \node[slab] at (0.66,0.34) {1};
  \node[slab] at (0.66,-0.36) {2};
\end{tikzpicture}
\figcap{7--29}
\end{minipage}\hfill
\begin{minipage}{0.48\linewidth}\centering
\begin{tikzpicture}[scale=1.124]
  \coordinate (A) at (0,2.0);
  \coordinate (C) at (0,1.0);
  \coordinate (B) at (0.643,1.234);
  \coordinate (Q) at (0.5303,0.5435);
  \draw[given] (A) -- (0,0);
  \draw[given] (A) -- (1.672,0.008);
  \draw[tick] (-0.13,1.0) -- (0.13,1.0);
  \draw[tick] ($(B)+(-50:0.13)$) -- ($(B)+(130:0.13)$);
  \draw[aux] ($(Q)+(150:0.26)$) arc (150:205:0.26);
  \draw[aux] ($(Q)+(-25:0.26)$) arc (-25:30:0.26);
  \dt{Q}
  \node[vlab,above] at (A) {$A$};
  \node[vlab,left] at (C) {$C$};
  \node[vlab,above right] at (B) {$B$};
  \node[vlab,below] at (Q) {$Q$};
\end{tikzpicture}
\figcap{7--30}
\end{minipage}
\end{bkfigure}}

% ------------------------------------------------------- Fig 7-31, 7-32
\newcommand{\FIGVIITHIRTYONETHIRTYTWO}{\begin{bkfigure}[\linewidth]
\begin{minipage}{0.48\linewidth}\centering
\begin{tikzpicture}[scale=1.171]
  \coordinate (A) at (0,0);
  \coordinate (C) at (1.5,0);          % r  on the lower side
  \coordinate (Cp) at (2.1,0);         % r' on the lower side
  \coordinate (D) at (30:1.5);         % r  on the upper side
  \coordinate (Dp) at (30:2.1);        % r' on the upper side
  \coordinate (P) at (1.633,0.4373);   % DC' meets D'C on the bisector
  \draw[given] (A) -- (0:2.55);
  \draw[given] (A) -- (30:2.55);
  % the two compass arcs, radius r and r', each cutting BOTH sides of the
  % angle -- at C, D and at C', D' respectively.
  \draw[aux] ([shift=(0:1.5)]A) arc (0:30:1.5);
  \draw[aux] ([shift=(0:2.1)]A) arc (0:30:2.1);
  \draw[aux] (D) -- (Cp);
  \draw[aux] (Dp) -- (C);
  \dt{P}
  \node[vlab,left] at (A) {$A$};
  \node[vlab] at (1.149,1.010) {$D$};
  \node[vlab] at (1.669,1.310) {$D'$};
  \node[vlab,below] at (C) {$C$};
  \node[vlab,below] at (Cp) {$C'$};
  \node[vlab] at (1.227,0.329) {$P$};
\end{tikzpicture}
\figcap{7--31}
\end{minipage}\hfill
\begin{minipage}{0.48\linewidth}\centering
\begin{tikzpicture}[scale=1.076]
  % upper: angle CAD bisected
  \coordinate (A) at (0,1.35);
  \draw[given] (A) -- ++(30:2.7);
  \draw[given] (A) -- ++(15:2.7);
  \draw[given] (A) -- ++(0:2.7);
  \node[vlab,left] at (A) {$A$};
  \node[vlab,above right] at ($(A)+(30:2.7)$) {$C$};
  \node[vlab,below right] at ($(A)+(0:2.7)$) {$D$};
  \node[slab] at ($(A)+(22.5:1.90)$) {1};
  \node[slab] at ($(A)+(7.5:1.90)$) {2};
  % lower: angle EBF bisected
  \coordinate (B) at (0,-0.55);
  \draw[given] (B) -- ++(30:2.7);
  \draw[given] (B) -- ++(15:2.7);
  \draw[given] (B) -- ++(0:2.7);
  \node[vlab,left] at (B) {$B$};
  \node[vlab,above right] at ($(B)+(30:2.7)$) {$E$};
  \node[vlab,below right] at ($(B)+(0:2.7)$) {$F$};
  \node[slab] at ($(B)+(22.5:1.90)$) {3};
  \node[slab] at ($(B)+(7.5:1.90)$) {4};
\end{tikzpicture}
\figcap{7--32}
\end{minipage}
\end{bkfigure}}

% ------------------------------------------------------- Fig 7-33, 7-34
\newcommand{\FIGVIITHIRTYTHREETHIRTYFOUR}{\begin{bkfigure}[\linewidth]
\begin{minipage}{0.48\linewidth}\centering
\begin{tikzpicture}[scale=1.000]
  \coordinate (B) at (0,0); \coordinate (A) at (0,1.85);
  \coordinate (C) at (1.25,0); \coordinate (D) at (2.85,0);
  \draw[given] (B) -- (D);
  \draw[given] (A) -- (B);
  \draw[given] (A) -- (C);
  \draw[given] (A) -- (D);
  \draw[tick] (0,0.20) -- (0.20,0.20) -- (0.20,0);
  \node[vlab,above] at (A) {$A$};
  \node[vlab,below left] at (B) {$B$};
  \node[vlab,below] at (C) {$C$};
  \node[vlab,below right] at (D) {$D$};
\end{tikzpicture}
\figcap{7--33}
\end{minipage}\hfill
\begin{minipage}{0.48\linewidth}\centering
\begin{tikzpicture}[scale=1.305]
  \coordinate (A) at (0,1.35); \coordinate (B) at (2.9,1.1);
  \coordinate (C) at (0.1,0);  \coordinate (D) at (3.0,0.35);
  \coordinate (O) at (1.913,0.7123);
  \draw[given] (A) -- (B);
  \draw[given] (C) -- (D);
  \draw[given] (A) -- (D);
  \draw[given] (C) -- (B);
  \node[vlab,above left] at (A) {$A$};
  \node[vlab,above right] at (B) {$B$};
  \node[vlab,below left] at (C) {$C$};
  \node[vlab,right] at (D) {$D$};
  \node[vlab] at (1.99,0.48) {$O$};
\end{tikzpicture}
\figcap{7--34}
\end{minipage}
\end{bkfigure}}

% ------------------------------------------------------- Fig 7-35, 7-36
\newcommand{\FIGVIITHIRTYFIVETHIRTYSIX}{\begin{bkfigure}[\linewidth]
\begin{minipage}{0.48\linewidth}\centering
\begin{tikzpicture}[scale=1.026]
  \coordinate (A) at (0,0);
  % B, A, F are collinear in the printed figure (158 and -22 degrees), which
  % is what makes Review Ex. 18's conclusion true: AC bisects angle DAB and
  % AE bisects angle DAF, hence angle CAE is right.
  \draw[given] (A) -- (158:1.70);     % B
  \draw[given] (A) -- (100.5:1.75);   % C
  \draw[given] (A) -- (43:1.75);      % D
  \draw[given] (A) -- (10.5:1.75);    % E
  \draw[given] (A) -- (-22:1.75);     % F
  \draw[aux] (A) -- (190.5:0.95);     % extension of AE (the exercise's hint)
  \node[vlab] at (0.10,-0.30) {$A$};
  \node[vlab,left] at (158:1.70) {$B$};
  \node[vlab,above] at (100.5:1.75) {$C$};
  \node[vlab,above right] at (43:1.75) {$D$};
  \node[vlab,right] at (10.5:1.75) {$E$};
  \node[vlab,right] at (-22:1.75) {$F$};
\end{tikzpicture}
\figcap{7--35}
\end{minipage}\hfill
\begin{minipage}{0.48\linewidth}\centering
\begin{tikzpicture}[scale=1.026]
  \coordinate (A) at (0,0);
  % AC bisects angle DAB ; angle CAE is right
  \draw[given] (A) -- (180:1.55);     % B
  \draw[given] (A) -- (127.5:1.75);   % C
  \draw[given] (A) -- (75:1.75);      % D
  \draw[given] (A) -- (37.5:1.75);    % E
  \draw[given] (A) -- (0:1.55);       % F
  \node[vlab] at (0,-0.32) {$A$};
  \node[vlab,left] at (180:1.55) {$B$};
  \node[vlab,above] at (127.5:1.75) {$C$};
  \node[vlab,above] at (75:1.75) {$D$};
  \node[vlab,right] at (37.5:1.75) {$E$};
  \node[vlab,right] at (0:1.55) {$F$};
\end{tikzpicture}
\figcap{7--36}
\end{minipage}
\end{bkfigure}}


% ---- local macros (hoist into book preamble) ----
% Numbered remark head. Sections 7-4 (pp.122) and 7-5 (pp.123-124) print
% "Remark 1.", "Remark 2." ... as running heads; brumfiel.sty's {remark}
% emits a bare "Remark." and would drop the book's numbering.
% Same shape as {remark}, with the number restored.
\newenvironment{VIIremarkn}[1]%
  {\par\medskip\noindent\emph{Remark #1.}\ }{\par\medskip}
% No \VIIstarnote: ch07 has starred exercises (7-2 #3, 7-4 #4-7, 7-7 #3,
% 7-8 #8-9, Algebra Review #7-9) but the book prints NO star footnote on any
% of pp.116-132 -- verified against the source pages, incl. the foot of
% p.118 where the chapter's first "*3." falls. ch02 owns the book-wide first
% use of the convention note (\IIstarnote); ch05 and ch10 omit it for the
% same reason.
% ---- end local macros ----

\begin{document}
\setcounter{chapter}{7}

\chapter*{\raggedright\Huge USE OF THE\\CONGRUENCE THEOREMS}
\addcontentsline{toc}{chapter}{7.\ Use of the Congruence Theorems}
\markboth{USE OF THE CONGRUENCE THEOREMS}{}
\vspace{-1em}\noindent\rule{\linewidth}{0.6pt}\vspace{1em}
\figurekey

%========================================================= 7-1
\section*{7--1\quad INTRODUCTION}
\addcontentsline{toc}{section}{7--1\quad Introduction}

In this chapter we do two things. In the first part we prove several
inequalities involving sides and angles of a triangle, ending up with the
\emph{triangle inequality}. This enables us to prove that a line segment is
the shortest path between two points. In the second part we prove that some of
the constructions we made with ruler and compass in Chapter 1 are correct.

%========================================================= 7-2
\section*{7--2\quad EXTERIOR ANGLES}
\addcontentsline{toc}{section}{7--2\quad Exterior angles}

In this section we prove a theorem about triangles which Euclid found to be
very useful. We need a definition.

\begin{definition}
If an angle is adjacent and supplementary to an angle of a triangle, it is
called an \emph{exterior angle} of the triangle. The other two angles of the
triangle are called the \emph{opposite interior angles}. For example,
$\ang{DBC}$ in Fig.~7--1 is an exterior angle and $\ang{A}$ and $\ang{C}$ are
interior angles and opposite to $\ang{DBC}$.
\end{definition}

\FIGVIIONE

\exercisegroup{7--1}
\begin{multicols}{2}
\begin{exlist}
\item In $\tri{ABC}$, how many exterior angles are there having vertex $B$?
  How are these exterior angles related to $\ang{B}$?
\item How many exterior angles does a triangle have?
\item Do the words \emph{opposite}, \emph{interior}, and \emph{exterior}, used
  in Definition 7--1, refer to properties or to relations?
\end{exlist}
\end{multicols}

\begin{theorem}
An exterior angle of a triangle is greater than either opposite interior angle.
\end{theorem}

We give two separate proofs of this theorem. The first proof may be found in
Hilbert's work, the second in Euclid's \emph{Elements}.

\noindent\emph{First Proof.}\quad
We know that, for any two angles, the first angle is either $>$, or $<$, or
$\cong$ the second. Therefore, we need only show that the exterior angle
cannot be $\cong$ or $<$ an opposite interior angle, and it must follow that
the exterior angle is $>$ an opposite interior angle. Suppose that the triangle
is labeled as in Fig.~7--2.

\FIGVIITWOTHREE

We first show that the exterior angle cannot be congruent to $\ang{C}$. Choose
$D$, as in the figure, such that $BD \cong AC$, and suppose that
$\ang{DBC} \cong \ang{C}$. Then $\tri{ACB} \cong \tri{DBC}$ because of the
S.A.S. theorem. Hence $\ang{CBA} \cong \ang{BCD}$. Therefore, because of the
addition of angles postulate, $\ang{ACD} \cong \ang{ABD}$. But $\ang{ABD}$ is
a straight angle, and so $\ang{ACD}$ is a straight angle. This is impossible.
(Why?) Hence $\ang{DBC}$ is \emph{not} congruent to $\ang{C}$.

We now show that it is impossible that $\ang{DBC} < \ang{ACB}$. Assume that
$\ang{DBC} < \ang{ACB}$. With $C$ as a vertex, let $\ang{ECB}$ be the angle, as
drawn in Fig.~7--3, that is congruent to $\ang{DBC}$. (Why is there such an
angle?) Then the ray $CE$ is in the interior of $\ang{ACB}$. (Why?) Hence the
segment $AB$ contains a point $F$ in the ray $CE$. (This is obvious from the
figure, and it may be proved with what we know about betweenness and interior
of angles.) Now consider $\tri{FCB}$. $\ang{DBC}$ is an exterior angle of this
triangle which is congruent to the opposite interior angle $FCB$. This has
already been shown to be impossible. Hence we cannot have
$\ang{DBC} < \ang{ACB}$. The only possibility then is
$\ang{DBC} > \ang{ACB}$. This completes the proof.

\noindent\emph{Second Proof.}\quad
Using the same labeling for the triangle, let $G$ be the point on $BC$, such
that $BG \cong GC$. Let $H$ be the point on the line $AG$ that is on the other
side of $G$ from $A$, such that $AG \cong GH$. Then
$\tri{GCA} \cong \tri{GBH}$. (Why?) Hence $\ang{C} \cong \ang{GBH}$. The ray
$BH$ is interior to $\ang{DBC}$. (This is obvious from Fig.~7--4, but it needs
proof. It is interesting to note that Euclid neglects to prove this. See
Exercise 3 below.) Therefore, $\ang{DBC} > \ang{HBC}$, and so
$\ang{DBC} > \ang{C}$. (Why?)

\FIGVIIFOUR

\begin{remark}
In Chapter 8, when we study parallel lines, we shall be able to prove that an
exterior angle is congruent to the ``sum'' of the opposite interior angles, and
from this fact, Theorem 7--1 follows easily. What is interesting is that
Theorem 7--1 may be proved without any knowledge of properties of parallel
lines.
\end{remark}

\exercisegroup{7--2}
\begin{exlist}
\item What kind of angles are the exterior angles of a right triangle?
\item Use Theorem 7--1 to prove that the angles opposite the legs of a right
  triangle are acute angles.
\item[\stex 3.]\stepcounter{exlisti} Prove that the ray $BH$ is interior to $\ang{DBC}$.
  [\emph{Hint:} first show that $H$ is on the same side of line $BC$ as $D$.
  Since $G$ is on the same side of $BD$ as $C$, if $H$ were on the opposite
  side of $BD$ from $C$, then the segment $GH$ would contain a point of the
  line $BD$. This is absurd.]
\item Prove the following corollary to Theorem 7--1. There is at most one line
  perpendicular to a given line $l$ from a given point $P$. [\emph{Hint:}
  Suppose that there were two perpendiculars, as in Fig.~7--5.]
\exfig{\FIGVIIFIVE}
\end{exlist}

%========================================================= 7-3
\section*{7--3\quad INEQUALITIES IN TRIANGLES}
\addcontentsline{toc}{section}{7--3\quad Inequalities in triangles}

In this section we prove some theorems about the relative sizes of sides and
angles in a triangle. We arrive at a proof of a theorem that, in effect, says
that a line segment is the shortest path between two points.

\begin{theorem}
If, in $\tri{ABC}$, $AB > AC$, then $\ang{C} > \ang{B}$, and conversely, if
$\ang{C} > \ang{B}$, then $AB > AC$.
\end{theorem}

\begin{remark}
The meaning of this theorem can be stated briefly as ``the greater side is
opposite the greater angle.''
\end{remark}

\FIGVIISIX

\noindent\emph{Proof.}\quad
To prove the first part of the theorem, suppose that $AB > AC$, as shown in
Fig.~7--6. Choose $C'$ on $AB$, such that $AC' \cong AC$. Then by
Theorem 7--1, $\ang{AC'C} > \ang{B}$. But $\ang{AC'C} < \ang{C}$. Hence
$\ang{C} > \ang{B}$.

The converse follows readily from the fact that any two segments, are either
congruent to each other or one is greater than the other. The argument is given
below.

Given that $\ang{C} > \ang{B}$, we wish to show that $AB > AC$. Exactly one of
the relations $AB < AC$, $AB \cong AC$, or $AB > AC$ is true. If
$AB \cong AC$, then we would have $\ang{C} \cong \ang{B}$, which contradicts
the given fact that $\ang{C} > \ang{B}$. If $AB < AC$, then by the first part
of the theorem, we would have $\ang{C} < \ang{B}$, which also contradicts the
given fact that $\ang{C} > \ang{B}$. Hence the only possibility is that
$AB > AC$.

\exercisegroup{7--3}
\begin{exlist}
\item Which angle in Fig.~7--7 is the smallest? Which is the largest? The
  numbers are lengths of the sides.
\exfig{\FIGVIISEVEN}
\item If in $\tri{EFG}$, $\ang{E}$ has a measure of $61^\circ$, and $\ang{F}$ a
  measure of $54^\circ$, which side is the longer, $FG$ or $EG$? Can you say
  anything yet about the length of $EF$?
\item In $\tri{PQR}$, $S$ is a point of $PR$, such that $QS$ bisects
  $\ang{PQR}$ (see Fig.~7--8). Prove that $PQ > PS$. [\emph{Hint:} Prove first
  that $\ang{PSQ} > \ang{PQS}$.]
\exfig{\FIGVIIEIGHT}
\end{exlist}

We now come to a very important and familiar inequality for triangles, called
the \emph{triangle inequality}.

\begin{theorem}
In any $\tri{ABC}$, $\sg{AB} + \sg{BC} > \sg{AC}$.
\end{theorem}

\FIGVIININETEN

\noindent\emph{Proof.}\quad
If $AC$ is not the longest of the sides, the inequality is obvious. Therefore
we may confine our attention to the case in which $\sg{AC} > \sg{AB}$ and
$\sg{AC} > \sg{BC}$, as in Fig.~7--9.

\begin{steps}
\step{There is a point $B'$ on $AC$, such that $AB \cong AB'$.}{Why?}
\step{$\ang{1} \cong \ang{2}$.}{Why?}
\step{$\ang{3} > \ang{1}$.}{Why?}
\step{$\ang{2} > \ang{4}$.}{Why?}
\step{$\ang{3} > \ang{4}$.}{This follows from Statements 2, 3, and 4. Why?}
\step{$\sg{B'C} < \sg{BC}$.}{Why?}
\step{$\sg{AB'} + \sg{B'C} < \sg{AB} + \sg{BC}$.}{Statements 1 and 6 and
      properties of numbers.}
\step{$\sg{AC} < \sg{AB} + \sg{BC}$.}{$\sg{AC} = \sg{AB'} + \sg{B'C}$.}
\end{steps}

\begin{remark}
This theorem is sometimes loosely stated as ``a straight line is the shortest
distance between two points.'' In fact, this last statement is sometimes used
as a postulate, although incorrectly, since one must first have a notion of
distance or of length of a segment.
\end{remark}

Having proved the triangle inequality, we now see that the lines we have been
talking about all along are, after all, ``straight'' in the precise sense of
the triangle inequality. Perhaps this can be made clearer by defining
\emph{polygonal path} and \emph{length of a polygonal path}.

\begin{definition}
Let $P_1$, $P_2,\ldots, P_n$ be any set of $n$ points (Fig.~7--10). The set of
points on all the segments $P_1P_2$, $P_2P_3,\ldots, P_{n-1}P_n$ is called a
\emph{polygonal path}. The path is said to join, or connect, the
\emph{endpoints} $P_1$ and $P_n$. The points $P_1,\ldots, P_n$ are called the
\emph{vertices}. In particular, if $P_1 = P_n$, the polygonal path is called a
\emph{polygon} and designated as polygon $P_1P_2\ldots P_{n-1}$.
\end{definition}

\begin{definition}
The \emph{length of the polygonal path} $P_1\ldots P_n$ is defined to be
$\sg{P_1P_2} + \sg{P_2P_3} + \cdots + \sg{P_{n-1}P_n}$.
\end{definition}

\begin{corollary}
The length of a polygonal path is greater than, or equal to, the length of the
segment joining its ends.
\end{corollary}

\exercisegroup{7--4}
\begin{exlist}
\item Two sides of a triangle have lengths of 5 and 7 in. What can you say
  about the third side?
\item Can the numbers 6, 8, and 15 be the lengths of the three sides of a
  triangle?
\item Prove that the difference between the lengths of two sides of a triangle
  is less than the length of the third side.
\item[\stex 4.]\stepcounter{exlisti} Use the triangle inequality and Fig.~7--11 to show that if
  $\ang{C} > \ang{B}$, then $AB > AC$.
\exfig{\FIGVIIELEVEN}
\item[\stex 5.]\stepcounter{exlisti} In Fig.~7--12, $AC \cong A'C'$, $BC \cong B'C'$, and
  $\ang{ACB} > \ang{A'C'B'}$. Prove that $\sg{AB} > \sg{A'B'}$.
  [\emph{Hint:} Show how to find a point $Q$ on $AB$, such that
  $QB \cong QB'$, and use $\sg{AB} = \sg{AQ} + \sg{QB} = \sg{AQ} + \sg{QB'}$.]
\exfig{\FIGVIITWELVE}
\item[\stex 6.]\stepcounter{exlisti} Use the result of Exercise 5 to prove that if, in two
  triangles, $\tri{ABC}$ and $\tri{A'B'C'}$, $AC \cong A'C'$,
  $BC \cong B'C'$, and $AB > A'B'$, then $\ang{C} > \ang{C'}$. Show that the
  statements $\ang{C} \cong \ang{C'}$ and $\ang{C} < \ang{C'}$ are false.
\item[\stex 7.]\stepcounter{exlisti} If $D$ in Fig.~7--13 is a point inside $\tri{ABC}$, show that
  $\sg{AD} + \sg{DC} < \sg{AB} + \sg{BC}$. [\emph{Hint:} Consider first
  $\tri{ABE}$ and $\tri{DEC}$.] Show also that $\ang{ADC} > \ang{ABC}$.
\exfig{\FIGVIITHIRTEEN}
\item By using Theorem 7--2, prove that an equiangular triangle is equilateral.
\item If, in Fig.~7--14, $l \perp m$, and $l' \perp m$, prove that $l$ and $l'$
  do not intersect. [\emph{Hint:} Suppose they do.]
\exfig{\FIGVIIFOURTEEN}
\item Show that a polygonal path connecting points $A$ and $B$ could have four
  vertices and still have the same length as $AB$.
\end{exlist}

%========================================================= 7-4
\section*{7--4\quad CIRCLES}
\addcontentsline{toc}{section}{7--4\quad Circles}

There are two ways that we could define a circle. One is to use the idea of
congruent segments, and the other is to use the idea of distance. Because two
segments are congruent if and only if they have the same length, these two ways
are equivalent.

\begin{definition}
A circle is determined by a point $O$, called the \emph{center}, and a positive
number $r$. The \emph{circle} is the set of all points $A$ of the plane, such
that $\sg{OA} = r$. If $A$ is a point of the circle, then the segment $OA$ is
called a \emph{radius} (plural, \emph{radii}). If $A$ and $B$ are on the
circle, and the segment $AB$ contains $O$, then $AB$ is called a
\emph{diameter}.
\end{definition}

\begin{VIIremarkn}{1}
You will note that we have not said that a circle is a ``curve,'' because
nowhere in our geometry have we defined what ``curve'' might mean. The circle
is a set of points, and in the definition, we have told only how to decide
whether any point $A$ is, or is not, a point of the circle.
\end{VIIremarkn}

\begin{VIIremarkn}{2}
According to the definition, a radius is a segment, that is, a set of points.
It is common usage to also call the real number $r$ the radius. We will also
use this language, so that the circle can be described as the ``circle with
center $O$ and radius $r$.'' Similarly, if $AB$ is a diameter, the length
$\sg{AB} = 2r$ will be called the diameter. No confusion will arise, because it
will always be clear whether we are talking about a set of points or a number.
\end{VIIremarkn}

\begin{VIIremarkn}{3}
We have not mentioned the inside or outside of the circle as yet, because these
sets have not yet been defined.
\end{VIIremarkn}

\begin{definition}
If a circle (Fig.~7--15) has center $O$ and radius $r$, the \emph{interior} of
the circle is the set of all points $B$, such that $\sg{OB} < r$. The
\emph{exterior} of the circle is the set of all points $C$, such that
$\sg{OC} > r$.

A \emph{chord} is a segment determined by two points $E$ and $F$ on the circle;
the line through $E$ and $F$ is a \emph{secant}.

A \emph{central angle} of a circle is an angle whose vertex is at the center of
the circle. An \emph{inscribed angle} (such as $\ang{EFG}$) is an angle whose
vertex is on the circle, and having points of the circle on each of the two
sides of the angle.
\end{definition}

\FIGVIIFIFTEEN

\exercisegroup{7--5}
\begin{exlist}
\item A circle has center $A$ and radius 5. The points $B$, $C$, $D$, $E$, are
  such that $\sg{AB} = 4$, $\sg{AC} = 5$, $\sg{AD} = 17$, and $\sg{AE} = 5$.
  Where are these points? That is, in what sets are they?
\item If $AB$ is a diameter of a circle with center $C$, what is $\ang{ACB}$?
\item In Fig.~7--16, the circle has center $O$. What kind of angles are
  $\ang{A}$, $\ang{B}$, $\ang{C}$, $\ang{D}$, and $\ang{E}$?
\exfig{\FIGVIISIXTEEN}
\item Give a definition of circle in terms of congruence instead of distance.
\item If $A$ is a point of a circle and $O$ is the center, what can be said
  about a point $B$, such that $\sg{OB} \leqq \sg{OA}$?
\end{exlist}

%========================================================= 7-5
\section*{7--5\quad RULER AND COMPASS}
\addcontentsline{toc}{section}{7--5\quad Ruler and compass}

In the preceding section we defined a circle in terms of distance. When drawing
a circle on a piece of paper, we use the fact, familiar from experience, that
the distance between the two points of the compass stays fixed. In this section
we define those drawings that are called ``ruler and compass constructions.''
Do not think that these are the only constructions possible. They are only the
simplest kind. The ruler used in these constructions is just an unmarked
straightedge.

\FIGVIISEVENTEEN

\begin{definition}
To make a \emph{ruler and compass} (\RC) \emph{construction} is to draw a
picture of a set of points, using only a ruler and compass. The ruler can only
be used to draw a line through two points. The compass can only be used to
draw a circle with a given center and radius.
\end{definition}

\begin{VIIremarkn}{1}
\RC{} constructions are sets that can be illustrated with certain mechanical
devices. Circles and lines exist in our geometry independently of these
mechanical devices. Interest in \RC{} constructions originates in the
simplicity of the tools needed to make them.
\end{VIIremarkn}

\begin{VIIremarkn}{2}
To draw a line or circle in an \RC{} construction, two points are needed. We
may not slide the ruler, mark it, or guess the opening of the compass. We must
have two points already constructed in order to draw a line or circle.
\end{VIIremarkn}

\begin{VIIremarkn}{3}
The study of \RC{} constructions at this point is not really necessary for the
logical development of our geometry, but it does help us to visualize geometric
figures, and it demonstrates the usefulness of what has already been studied.
\end{VIIremarkn}

\begin{VIIremarkn}{4}
It will be clear that the knowledge of \RC{} constructions is quite useful to a
draftsman.
\end{VIIremarkn}

%========================================================= 7-6
\section*{7--6\quad HISTORICAL NOTE}
\addcontentsline{toc}{section}{7--6\quad Historical note}

The ancient Greeks were much attracted by \RC{} constructions and attempted to
solve all sorts of problems with these two tools. However, they encountered
several problems that they could not handle using only ruler and compass. This
does not mean that these problems are not solvable---only that a ruler and
compass are not enough. Three of the problems that puzzled the Greeks have had
a long and interesting history, which we cannot pursue here. They are:

\begin{enumerate}[label=(\alph*), leftmargin=2.2em, itemsep=3pt]
\item \emph{The trisection of an angle.} Given any $\ang{A}$, the problem is to
  construct angles $B$, $C$, and $D$, as in Fig.~7--18, such that
  $\ang{B} \cong \ang{C} \cong \ang{D}$.
\item \emph{The duplication of a cube.} Suppose that a cube has side of length
  $x$. The problem is to construct the side of a cube having twice the volume
  of the given cube.
\item \emph{Squaring the circle.} Given a circle, the problem is to construct a
  square with the same area as the circle (Fig.~7--19).
\end{enumerate}

\FIGVIIEIGHTEENNINETEEN

All these problems have been completely resolved in modern times, and
\emph{all are impossible of solution with ruler and compass alone}. This is a
fact beyond dispute. It does not mean that in the trisection problem, for
example, angles $B$, $C$, and $D$ do not exist. It means that one cannot draw
them using ruler and compass alone. When you think about this a moment, it may
not be so surprising. No one is surprised that it is impossible to build a
large building with only a tack hammer. A ruler and compass alone are like a
tack hammer in so far as these problems are concerned.

There are, however, many constructions possible with ruler and compass. We will
study some of them.

%========================================================= 7-7
\section*{7--7\quad COPYING ANGLES AND TRIANGLES}
\addcontentsline{toc}{section}{7--7\quad Copying angles and triangles}

$\ang{A}$ and a ray from point $B$ are given, as in Fig.~7--20. We wish to
construct an angle at $B$ congruent to $\ang{A}$. Of course, from the existence
postulate for angles, there are two such angles, as shown by the dotted lines
in the figure. When we draw the angle at $B$ congruent to $\ang{A}$, we say
that we have ``copied'' $\ang{A}$.

\FIGVIITWENTYTWENTYONE

\begin{construction}
Draw a circle of radius $r$, with center $A$. This circle will intersect the
sides of $\ang{A}$ in two points $D$ and $E$ (Fig.~7--21). (Why?) Draw a circle
of the same radius $r$ with center $B$. This circle will intersect the ray from
$B$ in a point $F$. Draw a circle, with center $F$ and radius $\sg{DE}$. On a
given side of the line $BF$, the two circles will intersect in a point $G$.
$\ang{GBF}$ is congruent to $\ang{A}$.
\end{construction}

You should recognize that there are two things to prove. Because the entire
construction is centered around finding a point in which our circles intersect,
we must first prove that such a point exists in our geometry. Of course, when
we draw the picture, it is clear that there should be such a point. The real
question, however, is whether we can take the assumptions we have made and
prove that there is such a point. After we establish the existence of this
point, we must prove further that the angles are congruent to each other. These
two proofs should be considered separately.

\noindent\emph{Proof.} (a)\quad
(The existence of a point of intersection of the circles.) By the existence
postulate for angles, there is exactly one ray $BH$, on a given side of $BF$,
such that $\ang{FBH} \cong \ang{A}$. The existence of the \emph{ray} is thus
assured. The existence of the points $D$, $E$, and $F$ follows from theorems of
measurement and congruence. There is a point $G$ on the ray $BH$, such that
$BG \cong BF \cong AE \cong AD$. Since $\ang{A} \cong \ang{B}$,
$\tri{ADE} \cong \tri{BGF}$ by the S.A.S. theorem. Therefore,
$GF \cong DE$. Hence, the point $G$ is on the circle with radius $\sg{DE}$ and
center $F$. Since $G$ is also on the circle with center $B$ and radius
$\sg{BF}$, it is therefore on both circles and is a point of intersection of
the two circles.

\noindent\emph{Proof.} (b)\quad
(The congruence of the angles.) We now consider any point $G$ on both circles.
$AD \cong AE \cong BG \cong BF$. Since both circles have the same radius,
$DE \cong GF$, and $\tri{AED} \cong \tri{BFG}$ by the S.S.S. theorem. Hence
$\ang{A} \cong \ang{B}$.

In a similar manner we can copy a triangle.

\begin{construction}
Given $\tri{ABC}$ (Fig.~7--22) and a ray from a point $A'$, use the previous
construction to copy $\ang{A}$. Then draw a circle, with center $A'$ and radius
$\sg{AB}$, which intersects the other side of $\ang{A'}$ at $B'$. Then draw a
circle, with center $A'$ and radius $\sg{AC}$, which intersects the given ray
in $C'$. Then $\tri{ABC} \cong \tri{A'B'C'}$.
\end{construction}

\noindent\emph{Proof.}\quad
The proof is left to the student.

\FIGVIITWENTYTWOTWENTYTHREE

\begin{construction}
Given $\tri{ABC}$ (Fig.~7--23) and a ray from a point $A'$, draw a circle, with
center $A'$ and radius $\sg{AC}$, which intersects the ray in $C'$. Then draw
circles, with centers at $A'$ and $C'$ and radii $\sg{AB}$ and $\sg{BC}$,
respectively. These two circles will intersect, on a given side of the ray, in
exactly one point $B'$. We will then have $\tri{ABC} \cong A'B'C'$.
\end{construction}

\noindent\emph{Proof.}\quad
The proof is left to the student. Remember, you must first prove that $B'$
exists and then prove $\tri{ABC} \cong \tri{A'B'C'}$.

%========================================================= 7-8
\section*{7--8\quad OTHER CONSTRUCTIONS}
\addcontentsline{toc}{section}{7--8\quad Other constructions}

It should now be clear that proving the validity of certain constructions
involves first the proof of the existence of certain points and lines, and
secondly the proof that the construction produces the desired result. In the
remainder of our work on constructions, we will concentrate our efforts on the
second part. The proof of existence of certain points can actually be quite
hard at times. For example, when we try to prove that the construction for the
perpendicular bisector of a segment is valid, we need a theorem, Theorem 7--4,
that is very difficult to prove at this time.

\begin{theorem}
If $AB$ is a segment, and circles are drawn, with centers at $A$ and $B$ and
radius $r > \frac{1}{2}\sg{AB}$, then these circles intersect in exactly two
points $P$ and $Q$, on opposite sides of the line through $A$ and $B$
(Fig.~7--24).
\end{theorem}

\begin{remark}
From what we know about measurement of length, we know that every line segment
$AB$ has exactly one midpoint $C$, that is, a point $C$ such that
$\sg{AC} = \sg{BC} = \frac{1}{2}\sg{AB}$. Likewise, it is easy to show that
every angle has a bisector, that is, from the vertex of the angle, there is a
ray that forms, with the sides of the given angle, two congruent angles.
\end{remark}

\FIGVIITWENTYFOURTWENTYFIVE

The constructions of this section show how to make \RC{} constructions for
midpoints of segments and bisectors of angles; but the \emph{existence} of
these midpoints and bisectors has been obtained independently of these
constructions.

\begin{construction}
Draw circles, with centers $A$ and $B$ and equal radii greater than
$\frac{1}{2}\sg{AB}$. These circles intersect at points $P$ and $Q$. The line
$PQ$ is the desired perpendicular bisector (Fig.~7--25).
\end{construction}

\noindent\emph{Proof.}
\begin{steps}
\step{$AP \cong PB \cong AQ \cong QB$.}{Construction.}
\step{$\tri{APQ} \cong \tri{BPQ}$.}{S.S.S. theorem.}
\step{$PQ$ intersects line $AB$ in $C$.}{$P$ and $Q$ are on opposite sides of
      $AB$.}
\step{$\tri{APC} \cong \tri{BPC}$.}{S.A.S. theorem.}
\step{$AC \cong CB$, and $\ang{ACP} \cong \ang{BCP}$.}{Statement 4.}
\step{$PQ$ is the $\perp$ bis $AB$.}{Statement 5.}
\end{steps}

\exercisegroup{7--6}
\begin{multicols}{2}
\begin{exlist}
\item Prove that each segment has exactly one perpendicular bisector.
\item Prove that every point on the perpendicular bisector of $AB$ is
  equidistant from $A$ and $B$.
\item Prove that if a point $R$ is equidistant from $A$ and $B$, then it lies
  on the perpendicular bisector of $AB$.
\item Construct the perpendicular bisector of a segment by constructing two
  points on the same side of the line and equidistant from the ends of the
  segment.
\end{exlist}
\end{multicols}

\begin{construction}
To construct the perpendicular to a line at a point $P$ on the line
(Fig.~7--26), draw a circle, with center $P$ and arbitrary radius, intersecting
the line in points $A$ and $B$. Then construct the perpendicular bisector of
$AB$.
\end{construction}

\begin{construction}
To construct a perpendicular to a line from a point not on the line, choose a
point $A$ on the line, and draw the circle with center $P$ and radius $AP$
(Fig.~7--27). The circle will intersect the line in a second point $B$ (if $A$
is properly chosen). Then construct the perpendicular bisector of $AB$, which
passes through $P$.
\end{construction}

\FIGVIITWENTYSIXSEVENEIGHT

\noindent\emph{Proof.}\quad
We first show that there actually is a line through $P$ perpendicular to the
line. Then we show that the \RC{} construction produces this line.

First choose any point $A$ on $l$ and any other point $C$ on $l$ (Fig.~7--28).
Let $AQ$ be the ray, on the opposite side of $l$ from $P$, such that
$\ang{PAC} \cong \ang{QAC}$. (There is only one such ray, by the existence
postulate for angles.) Choose $P'$ on $AQ$, such that $AP' \cong AP$. Then
$PP'$ intersects $l$ in a point $D$. (Why?) If $A \neq D$, then triangles are
formed, and by S.A.S., $\tri{APD} \cong \tri{AP'D}$. Therefore, $\ang{ADP}$ and
$\ang{ADP'}$ are right angles, and $PP' \perp l$. If $A = D$, there are no
triangles, but then, because of the way $AQ$ was selected, $PP' \perp l$.
(Why?)

Having established that a perpendicular from $P$ exists, we now show that the
\RC{} construction gives us this perpendicular. If $D \neq A$, there is a point
$B$ on the opposite side of $D$ from $A$, so that $AD \cong BD$. It follows
that the circle with center $P$ and radius $AP$ does intersect the line in a
second point $B$, and this was the main part to be proved in the construction.
Now the perpendicular bisector of $AB$ can be drawn, and $P$ lies on it.
(Why?)

\exercisegroup{7--7}
\begin{exlist}
\item Construct a right triangle, having been given the lengths of the two
  legs.
\item Construct points $D$, $E$, and $F$ on the segment $AB$, such that
  $\sg{AD} = \sg{DE} = \sg{EF} = \sg{FB} = \frac{1}{4}\sg{AB}$.
\item[\stex 3.]\stepcounter{exlisti} Suppose that $P$ is not on a line $l$, and that a given circle
  with center at $P$ intersects $l$ in exactly one point $Q$. Prove that $PQ$
  is perpendicular to $l$. [\emph{Hint:} If $PQ$ were not perpendicular, then
  the construction in the proof of Construction 7--6 would lead to a
  contradiction.]
\item Construct the perpendicular from $P$ by choosing $A$, as in Fig.~7--29,
  constructing $\ang{2} \cong \ang{1}$, and point $P'$, such that
  $AP \cong AP'$.
\exfig{\FIGVIITWENTYNINE}
\end{exlist}

\begin{construction}
To bisect $\ang{A}$, draw a circle, with center $A$, intersecting the sides of
$\ang{A}$ in points $B$ and $C$ (Fig.~7--30). Then construct a point $Q$ on the
perpendicular bisector of $BC$ and in the interior of $\ang{A}$. The ray $AQ$
bisects $\ang{A}$.
\end{construction}

\noindent\emph{Proof.}\quad
The point $Q$ can be found by a previous construction.
$\tri{ABQ} \cong \tri{ACQ}$, by S.S.S. Hence $\ang{BAQ} \cong \ang{CAQ}$.

\FIGVIITHIRTY

\exercisegroup{7--8}
\begin{exlist}
\item Draw an acute angle, and construct its bisector.
\item Construct and bisect a right angle.
\item Bisect an obtuse angle.
\item Draw a triangle, and construct the bisectors of its angles. (They should
  meet in one point.)
\item If $P$ is on the bisector of $\ang{A}$, and $B$ and $C$ are points on the
  sides of $\ang{A}$, such that $AB = AC$, prove that $PB \cong PC$.
\item Prove that an angle has only one bisector.
\item Suppose that $B$ and $C$ are points on the two sides of $\ang{A}$, such
  that $AB \cong AC$. Show that if $M$ is the midpoint of $BC$, then $AM$
  bisects $\ang{A}$.
\item[\stex 8.]\stepcounter{exlisti} Figure 7--31 shows another construction for the angle bisector.
  Consider two circles, with center $A$ and radii $r$ and $r'$, such that
  $r' > r$. These circles intersect the sides of $A$ in four points $C$, $C'$,
  $D$, and $D'$. The lines through these points intersect in a point $P$ on the
  bisector. (Why?)
\exfig{\FIGVIITHIRTYONE}
\item[\stex 9.]\stepcounter{exlisti} If, in Fig.~7--32, $\ang{CAD} \cong \ang{EBF}$,
  $\ang{1} \cong \ang{2}$, and $\ang{3} \cong \ang{4}$, then
  $\ang{2} \cong \ang{4}$. \emph{Note:} This theorem is often stated as
  ``halves of congruent angles are congruent to each other.'' Prove this
  theorem.
\exfig{\FIGVIITHIRTYTWO}
\end{exlist}

%========================================================= Review
\section*{REVIEW OF CHAPTER 7}
\addcontentsline{toc}{section}{Review of Chapter 7}

In this chapter we have used the congruence theorems to study inequalities in
triangles and verify constructions.

\begin{enumerate}[label=(\alph*), leftmargin=2.2em, itemsep=3pt]
\item Make a list of all the theorems on inequalities in triangles, then study
  their proofs.
\item Make a list of the different \RC{} constructions. Prove them correct
  without looking at the book.
\end{enumerate}

Review the following definitions, which you met for the first time in this
chapter: exterior angle, circle, interior of circle, exterior of circle,
radius, diameter, chord, secant, central angle, inscribed angle, \RC{}
construction.

\begin{center}\bfseries Review Exercises\end{center}\nopagebreak
\begin{exlist}
\item Construct an angle of $45^\circ$.
\item Construct an angle of $60^\circ$, of $30^\circ$.
\item Construct an angle of $15^\circ$, of $75^\circ$.
\item Construct an equilateral triangle.
\item Construct a right triangle having acute angles of $60^\circ$ and
  $30^\circ$.
\item Draw an obtuse triangle, and construct the bisectors of its angles.
\item Draw an angle. Copy it, using ruler and compass.
\item Draw three separate segments. Construct a triangle having sides congruent
  to the given segments.
\item Draw a line, and mark a point $P$ not on the line. Construct the
  perpendicular to the line, passing through $P$. How do you know that there is
  only this one perpendicular?
\item Draw a triangle. Construct the altitudes of the triangle.
\item Draw a triangle. Construct the perpendicular bisectors of its sides.
\item Construct an isosceles triangle if the vertex angle and the congruent
  sides are given.
\item Draw a segment. Construct a circle having this segment as a diameter.
\item In a right triangle, which side is longest? Why?
\item Prove that the sum of the lengths of three sides of a quadrilateral is
  greater than the length of the fourth side.
\item $\ang{ABC}$ is a right angle, and $C$ is between $B$ and $D$. Prove that
  $AD > AC$. (Fig.~7--33)
\exfig{\FIGVIITHIRTYTHREE}
\item Prove that $\sg{AD} + \sg{CB} > \sg{AB} + \sg{CD}$. (Fig.~7--34)
\exfig{\FIGVIITHIRTYFOUR}
\item $AE$ bisects $\ang{DAF}$, and $AC$ bisects $\ang{DAB}$. (Fig.~7--35)
  Prove that $\ang{CAE}$ is right. [\emph{Hint:} Extend $AE$.]
\exfig{\FIGVIITHIRTYFIVE}
\item $AC$ bisects $\ang{DAB}$, and $\ang{CAE}$ is right. Prove that $AE$
  bisects $\ang{DAF}$. (Fig.~7--36)
\exfig{\FIGVIITHIRTYSIX}
\end{exlist}

\begin{center}\bfseries ALGEBRA REVIEW\end{center}\nopagebreak
\begin{multicols}{2}
\begin{exlist}
\item In $\tri{ABC}$, $AB < BC < CA$. The greatest angle in this triangle has a
  measure of $32^\circ$ more than the measure of the smallest angle. The sum of
  the measures of these two angles is $120^\circ$. Determine the number of
  degrees in the smallest angle in this triangle.
\item An exterior angle at $A$ of $\tri{ABC}$ has a measure of $65^\circ$. The
  sum of the measures of angles $A$ and $B$ is $145^\circ$, and the sum of the
  measures of angles $A$ and $C$ is $150^\circ$. Determine the number of
  degrees in each angle of the triangle.
\item Points $O$, $A$, and $B$ are collinear. $\sg{OA}$ is 7 in. more than
  $\sg{BA}$. $A$ is not between $O$ and $B$. Determine the length of $OB$.
\item $\sg{AB} = 17$, and $\sg{AC} = 12$. The number $\sg{BC}$ is a root of the
  equation $(x-6)(x-4) = 0$. Determine $\sg{BC}$.
\item If $B$ is between $A$ and $C$, $\sg{AC}/\sg{BC} = \frac{5}{3}$, and
  $\sg{AB} = 10$ in., determine the lengths of $AC$ and $BC$.
\item $B$ is between $A$ and $C$, and $\sg{AB}/\sg{AC} = \frac{3}{5}$. The
  midpoint of $AC$ is $M$, and $\sg{MB} = 3$. Determine $\sg{AC}$.
\item[\stex 7.]\stepcounter{exlisti} $A$, $B$, $C$ and $D$ are points in this order on a line. The
  length of $AC$ is twice the length of $BD$. The length of $AB$ is four times
  the length of $CD$. The length of $BC$ is 4. Determine the lengths of $AB$
  and $AD$.
\item[\stex 8.]\stepcounter{exlisti} The difference of the squares of two numbers is 33. The sum of
  these numbers is 11. Determine the numbers.
\item[\stex 9.]\stepcounter{exlisti} Analyze the following mathematical trick. Choose any number;
  multiply this number by 2; add 4 to this product; multiply by 5; divide by
  10; subtract from this result the original number. The result will be 2, no
  matter what number was originally chosen.
\end{exlist}
\end{multicols}

\end{document}
